\section{Background} \label{sec:background}
\paragraph{Group Fair Classification}
In fair classification, we consider labeled examples $x, a, y \sim p_{\text{data}}$ where $y \in \mathcal{Y}$ are labels we wish to predict, $a \in \mathcal{A}$ are \emph{sensitive} attributes, and $x \in \mathcal{X}$ are non-sensitive attributes. 
The goal is to learn a classifier $\hat y = g(x, a)$ (or $\hat y = g(x)$) which is predictive of $y$ and achieves certain group fairness criteria w.r.t. $a$.
These criteria are typically written as independence properties of the various random variables involved.
In this paper we focus on demographic parity, which is satisfied when the predictions are independent of the sensitive attributes: $\hat y \perp a$.
It is often impossible or undesirable to satisfy demographic parity exactly (i.e. achieve complete independence).
In this case, a useful metric is \textit{demographic parity distance}:
\begin{equation}
    \Delta_{DP} = | \E [\bar y = 1 | a = 1] - \E [\bar y = 1 | a = 0] |
\end{equation}
where $\bar y$ here is a binary prediction derived from model output $\hat y$.
When $\Delta_{DP}=0$, demographic parity is achieved; in general, lower $\Delta_{DP}$ implies less unfairness.

Our work differs from the fair classification setup as follows:
We consider several sensitive attributes at once, and seek fair outcomes with respect to each; individually as well as jointly (cf. subgroup fair classification, \citet{kearns2017preventing,johnson2018multicalibration});
also, we focus on representation learning rather than classification, with the aim of enabling a range of fair classification tasks downstream.

\paragraph{Fair Representation Learning}
In order to flexibly deal with many label and sensitive attribute sets, we employ representation learning to compute a compact but predicatively useful encoding of the dataset that can be flexibly adapted to different fair classification tasks.
As an example, if we learn the function $f$ that achieves independence in the representations $z \perp a$ with $z=f(x, a)$ or $z=f(x)$, then any predictor derived from this representation will also achieve the desired demographic parity, $\hat y \perp a$ with $\hat y = g(z)$.

The fairness literature typically considers binary labels and sensitive attributes: $\mathcal{A} = \mathcal{Y} = \{0, 1\}$.
In this case, approaches like regularization \citep{zemel2013learning} and adversarial regularization \citep{edwards2015censoring, madras2018learning} are straightforward to implement. 
We want to address the case where $a$ is a vector with many dimensions.
Group fairness must be achieved for each of the dimensions in $a$ (age, race, gender, etc.) and their combinations.

\paragraph{VAE}
The vanilla Variational Autoencoder (VAE) \citep{kingma2013auto} is typically implemented with an isotropic Gaussian prior $p(z) = \mathcal{N}(0, I)$.
The objective to be maximized is the Evidence Lower Bound (a.k.a., the ELBO), 
\begin{align}
    L_{\text{VAE}}(p,q) = \E_{q(z|x)} \left[ \log p(x|z) \right] - D_{KL} \left[ q(z|x) || p(z)  \right], \nonumber
\end{align}
which bounds the data log likelihood $\log p(x)$ from below for any choice of $q$.
The encoder and decoder are often implemented as Gaussians 
\begin{align}
    q(z|x) &= \mathcal{N}(z|\mu_q(x),  \Sigma_q(x)) \nonumber \\
    p(x|z) &= \mathcal{N}(x|\mu_p(z), \Sigma_p(z)) \nonumber
\end{align}
whose distributional parameters are the outputs of neural networks $\mu_q(\cdot)$, $\Sigma_q(\cdot)$, $\mu_p(\cdot)$, $\Sigma_p(\cdot)$, with the $\Sigma$ typically exhibiting diagonal structure.
For modeling binary-valued pixels, a Bernoulli decoder $p(x|z)=\text{Bernolli}(x|\theta_p(z))$ can be used.
The goal is to maximize $L_{\text{VAE}}$---which is made differentiable by reparameterizing samples from $q(z|x)$---w.r.t. the network parameters.

\paragraph{$\beta$-VAE}
\citet{higgins2016beta} modify the VAE objective:
\begin{align}
    L_{\beta\text{VAE}}(p,q) = \E_{q(z|x)} \left[ \log p(x|z) \right] - \beta D_{KL} \left[ q(z|x) || p(z) \right ]. \nonumber
\end{align}
The hyperparameter $\beta$ allows the practitioner to encourage the variational distribution $q(z|x)$ to reduce its KL-divergence to the isotropic Gaussian prior $p(z)$.
With $\beta>1$ this objective is a valid lower bound on the data likelihood. 
This gives greater control over the model's adherence to the prior.
Because the prior factorizes per dimension $p(z) = \prod_j p(z_j)$, \citet{higgins2016beta} argue that increasing $\beta$ yields ``disentangled'' latent codes in the encoder distribution $q(z|x)$.
Broadly speaking, each dimension of a properly disentangled latent code should capture no more than one semantically meaningful factor of variation in the data.
This allows the factors to be manipulated in isolation by altering the per-dimension values of the latent code.
Disentangled autoencoders are often evaluated by their sample quality in the data domain, but we instead emphasize the role of the encoder as a representation learner to be evaluated on downstream fair classification tasks.

\paragraph{FactorVAE and $\beta$-TCVAE}
\citet{kim2018disentangling} propose a different variant of the VAE objective:
\begin{align}%\label{eq:factorvae}
    L_{\text{FactorVAE}}(p,q) = L_{\text{VAE}}(p,q) - \gamma D_{KL}(q(z)||\prod_{j} q(z_j)). \nonumber
\end{align}
The main idea is to encourage factorization of the aggregate posterior $q(z)=\E_{p^{\text{data}}(x)} \left[ q(z|x) \right]$ so that $z_i$ correlates with $z_j$ if and only if $i=j$.
The authors propose a simple trick to generate samples from the aggregate posterior $q(z)$ and its marginals $\{q(z_j)\}$ using shuffled minibatch indices, then approximate the $D_{KL}(q(z)||\prod_{j} q(z_j))$ term using the cross entropy loss of a classifier that distinguishes between the two sets of samples, which yields a mini-max optimization.

\citet{chen2018isolating} show that the $D_{KL}(q(z)||\prod_{j} q(z_j))$ term above---a.k.a. the ``total correlation'' of the latent code---can be naturally derived by decomposing the expected KL divergence from the variational posterior to prior:
\begin{align} \nonumber% \label{eq:kl-decomposition}
    \E_{p^{\text{data}}(x)}[&D_{KL}(q(z|x)||p(z))] = \\
    &\quad\quad D_{KL}(q(z|x)p^{\text{data}}(x)||q(z)p^{\text{data}}(x)) \nonumber \\
        &\quad\quad+ D_{KL}(q(z)||\prod_j q(z_j))
     \nonumber\\
    &\quad\quad + \sum_j D_{KL} \left[ q(z_j) || p(z_j) \right]. \nonumber
\end{align}
They then augment the decomposed ELBO to arrive at the same objective as \citet{kim2018disentangling}, but optimize using a biased estimate of the marginal probabilities $q(z_j)$ rather than with the adversarial bound on the KL between aggregate posterior and its marginals.
